% ---- pg-509 （书页 497） ----

但是我们仍然保留性别进行分类. 我们将鲑鱼 $j$ 分类为 1(阿拉斯加), 若线性分类函数
\[
\boldsymbol{\beta}'\mathbf{z}=3.5054+0.2816\ \text{性别}+0.1264\ \text{淡水生长}+0.0486\ \text{海水生长}\ge 0
\]
错误分类的观测见下表

\begin{mvtable}
\begin{tabular}{cccccc}
\toprule
行 & 总体 & 性别 & 淡水 & 海水 & 线性分类函数 \\
\midrule
2 & 1 & 1 & 131 & 355 & 3.093 \\
12 & 1 & 2 & 123 & 372 & 1.537 \\
13 & 1 & 1 & 123 & 372 & 1.255 \\
30 & 1 & 2 & 118 & 381 & 0.467 \\
51 & 2 & 1 & 129 & 420 & $-0.319$ \\
68 & 2 & 2 & 136 & 438 & $-0.028$ \\
71 & 2 & 2 & 90 & 385 & $-3.266$ \\
\bottomrule
\end{tabular}
\end{mvtable}

根据错误分类数据得到混淆矩阵:

\begin{mvtable}
\begin{tabular}{c|l|c|c|}
\multicolumn{2}{c}{} & \multicolumn{2}{c}{预测类属}\\
\multicolumn{2}{c}{} & \multicolumn{1}{c}{$\pi_1$: 阿拉斯加} & \multicolumn{1}{c}{$\pi_2$: 加拿大}\\ \cline{3-4}
\multirow{2}{*}{实际类属} & $\pi_1$: 阿拉斯加 & 46 & 4 \\ \cline{3-4}
 & $\pi_2$: 加拿大 & 3 & 47 \\ \cline{3-4}
\end{tabular}
\end{mvtable}

近似的误差率用百分比表示为
\[
\mathrm{APER}=\frac{4+3}{50+50}\times 100=7\%
\]
\hfill$\blacksquare$

在作逻辑斯蒂分类时, 最好使用折刀法(保留法)对误分类概率进行估计. 但是这种方法在通用的主要统计软件包中还没有.

\medskip
\noindent\textbf{显示板 11.2}\quad 利用 PROC LOGISTIC 对鲑鱼数据的 SAS 分析

\begin{mvbox}
{\ttfamily\small
title 'Logistic Regression and Discrimination';\hfill\textbf{\rmfamily PROGRAM COMMANDS}\\
data salmon;\\
infile 'T11-2.dat';\\
input country gender freshwater marine;\\
proc logistic desc;\\
model country = gender freshwater marine / expb;
\par}
\end{mvbox}

\begin{mvbox}
{\ttfamily\small
\hfill\textbf{OUTPUT}\\
\par\vspace{6pt}
{\centering\small Logistic Regression and Discrimination\par}
\vspace{4pt}
{\centering\small The LOGISTIC procedure\par}
\vspace{8pt}
{\centering\small Model Information\par}
\vspace{4pt}
{\centering\small
\begin{tabular}{ll}
Model & binary logit\\
\end{tabular}\par}
\vspace{8pt}
{\centering\small Response Profile\par}
\vspace{4pt}
{\centering\small
\begin{tabular}{lcc}
Ordered Value & country & \makecell{Total\\Frequency}\\
1 & 2 & 50\\
2 & 1 & 50\\
\end{tabular}\par}
\par}
\end{mvbox}
